\documentclass[a4paper]{article}
%\usepackage[latin9]{inputenc}
\usepackage{amsmath}
\usepackage{amssymb}
\usepackage{graphicx}
\usepackage{mdframed}
%\usepackage{esint}
%\usepackage{cancel}
%\usepackage{float}
%\usepackage{graphicx}
%\usepackage{caption}
%\usepackage{subcaption}
%\usepackage{subfig}
%\usepackage{bbold}
%\usepackage{amsmath}
%\usepackage{amssymb}
%\usepackage{slashed}
%\usepackage{breqn}
%\usepackage{feynmf}

\makeatletter
%%%%%%%%%%%%%%%%%%%%%%%%%%%%%% User specified LaTeX commands.

\title{Superposition}
\author{Reuven Balkin}

\makeatother
\setlength{\textwidth}{500pt}
\setlength{\voffset}{-75pt}
\setlength{\textheight}{700pt}
\setlength{\hoffset}{-75pt}
\begin{document}

\maketitle

\makeatother
\noindent
\newcommand{\dif}{\mathtext{d}}
\section{Linear model}
Our setup follows a toy model presented in 'Toy Models of Superposition', where we attempt to reconstruct a vector $x \in \mathbb{R}^n$ by first projecting him to a lower-dimenstional space $\mathbb{R}^m$. We define
\begin{align}
h &= w x \,,
\\
x &= w^T h + b\,,
\\
x' &=  w^T w x + b\,.
\end{align}
The problem does not have a preferred basis (for any solution $w$, $Ow$ is also a solution where $O$ is an orthogonal rotation, therefore the sensible object to study is $w^T w$.
We start with the simplest problem $n=2,m=1$. We parameterize
\begin{align}
w &= (c_\theta, s_\theta)\,,
\\
w^T w &= \begin{pmatrix}
c^2_\theta & \frac12 s_{2\theta} 
\\
\frac12 s_{2\theta} & s^2_\theta
\end{pmatrix}\,.
\end{align}
The network is trained on the set $x_{Ji}$ where $J=1,....,k$ and $i=1,2$. Each $x_{Ji}$ is an i.i.d with probability $1-p_i$ to be 0, and probability $p_i$ to be distributed uniformly between 0 and 1.  The loss function is a weighted MSE, which can be written as
\begin{align}
\mathcal{L} &= \frac{1}{2k} \sum_{J}\sum_{i} I_i (x_{Ji}-x'(x_{jI}))^2 \equiv   \frac{1}{2m} \sum_{J} \mathcal{L}_J
%\\
%&= \frac{1}{2m} \sum_{J}\sum_{i} I_i \left(F_{i,\theta}x_{Ji}-\frac12 s_{2\theta} x_{J\bar{i}}-b_i \right)^2
%\\
%F_{i,\theta} &= \begin{cases}
%s^2_\theta \;\;\; &i=1
%\\
%c^2_\theta & i=2
%\end{cases}\,, \;\;\;\; \bar{i} \equiv (i+1) \mod 2
\end{align}
where $I_i $ are the weights for each of the features. For $k \gg 1$, the loss function is a narrowly-distributed variable, and we can consider its mean in order to study the phases of the system. The fundamental i.i.d variables satify
\begin{align}
E[x_{Ji}] &= \frac{p_i}{2}\,, \;\;\;\;\;
E[x^2_{Ji}] = \frac{p_i}{3}\,.
\end{align}
We find
\begin{align}
E[\mathcal{L}] &= \frac{1}{6}I_1 s_\theta^4 p_1 + \frac{1}{6}I_2 c_\theta^4 p_2 
+ \frac{s_{2\theta}^2}{24}\sum_i I_i p_{\bar{i}} 
+ \frac{1}{2}\sum_i I_i b_i^2 \nonumber\\
&\quad - \frac{s_{2\theta}}{8}\left(I_1 s_\theta^2 + I_2 c_\theta^2\right)p_1 p_2 
- \frac{1}{2}\left(I_1 s_\theta^2 b_1 p_1 + I_2 c_\theta^2 b_2 p_2\right) 
+ \frac{s_{2\theta}}{4}\sum_i I_i p_{\bar i} b_i
\end{align}
We find the solutions for $b$ in the minimum to be,
\begin{align}
b_1^* &= \frac{s_\theta}{2}(p_1 s_\theta - p_2 c_\theta)\\
b_2^* &= \frac{c_\theta}{2}(p_2 c_\theta - p_1 s_\theta)\,,
\end{align}
note that $\frac{\partial^2 E[\mathcal{L}]}{\partial b_i^2}=I_i>0$. Therefore the loss function become a function of $\theta$ alone, namely
\begin{align}
E[\mathcal{L}]\big|_{b^*} &= \frac{1}{6}(I_1 p_1 s_\theta^4 + I_2 p_2 c_\theta^4) + \frac{s_{2\theta}^2}{24}(I_1 p_2 + I_2 p_1) - \frac{(I_1 s_\theta^2 + I_2 c_\theta^2)(p_1^2 s_\theta^2 + p_2^2 c_\theta^2)}{8}\,.
\end{align}
and the first and second derivatives
\begin{align}
\frac{\partial E[\mathcal{L}]}{\partial\theta}\Bigg|_{b^*} &= \frac{s_{2\theta}}{6}(I_1 p_1 - I_2 p_2) - \frac{s_{2\theta}}{4}(I_1 p_1^2 s_\theta^2 - I_2 p_2^2 c_\theta^2)
\\
\frac{\partial^2 E[\mathcal{L}]}{\partial\theta^2}\Bigg|_{b^*} &= \frac{2c_{2\theta}}{6}(I_1 p_1 - I_2 p_2) - \frac{2c_{2\theta}}{4}(I_1 p_1^2 s_\theta^2 - I_2 p_2^2 c_\theta^2) - \frac{s_{2\theta}^2}{4}(I_1 p_1^2 + I_2 p_2^2)
\end{align}
Let us analyze the possible stationery points. First we have trivial solution at $\theta=0,\pi/2$, for which
\begin{align}
\frac{\partial^2 E[\mathcal{L}]}{\partial\theta^2}\Bigg|_{b^*,\theta=0} &= \frac{1}{3}(I_1 p_1 - I_2 p_2) +\frac{1}{2}I_2 p_2^2
\\
\frac{\partial^2 E[\mathcal{L}]}{\partial\theta^2}\Bigg|_{b^*,\theta=\pi/2} &= -\frac{1}{3}(I_1 p_1 - I_2 p_2) +\frac{1}{2}I_1 p_1^2  
\end{align}
Thus, we have two clear cases where if
\begin{align}
I_1p_1- I_2 p_2 < -(3/2) I_2 p_2^2 \;\; &\to \;\; \theta=\pi/2\,,
\\
I_1p_1 - I_2 p_2 > (3/2) I_1 p_1^2 \;\; &\to \;\; \theta=0\,,
\end{align}
For $-(3/2) I_2 p_2^2<I_1p_1 - I_2 p_2 < (3/2) I_1 p_1^2$ both are minima, with one being the deeper the two. We note however that this region becomes tighter with scarcity, since the interval edges scale like $\mathcal{O}(p^2)$. Indeed, the problem becomes much simpler if we neglect $\mathcal{O}(p^2)$ terms, in which case $\text{max}[I_1 p_1,I_2 p_2]$ selects which feature is represented. 

There is an additional non-trivial solution when $s_{2\theta}\neq0$,
%\begin{align}
%I_1 p_1 -\frac{3}{2}I_1 p_1^2 s_\theta^2  = I_2 p_2 - \frac{3}{2} I_2 p_2^2 c_\theta^2
%\end{align}
\begin{align}
c_{\theta_*}^2 = \frac{I_1 p_1^2 -\frac{2}{3}(I_1 p_1 - I_2 p_2) }{ I_2 p_2^2 +I_1 p_1^2}
\end{align}
The second derivative at $\theta_*$ is given by
\begin{align}
\frac{\partial^2 E[\mathcal{L}]}{\partial\theta^2}\Bigg|_{b^*,\theta=\theta_*} = -\frac{s^2_{2\theta}}{4}(I_1 p_1^2+I_2 p_2^2)<0 \,.
\end{align}
Therefore the only minima are $0$ and $\pi/2$, and superposition does not occur.

\subsection*{Phase boundary}
To determine which of the two minima is the global one, we evaluate $E[\mathcal{L}]\big|_{b^*}$ at each. At $\theta=0$ the network represents only feature~1 ($w^Tw = \mathrm{diag}(1,0)$), so the loss reduces entirely to the residual on feature~2:
\begin{align}
E[\mathcal{L}]\big|_{b^*,\,\theta=0} &= I_2 p_2\!\left(\frac{1}{6} - \frac{p_2}{8}\right).
\end{align}
Similarly, at $\theta=\pi/2$ only feature~2 is represented and
\begin{align}
E[\mathcal{L}]\big|_{b^*,\,\theta=\pi/2} &= I_1 p_1\!\left(\frac{1}{6} - \frac{p_1}{8}\right).
\end{align}
Feature~1 is represented ($\theta=0$ is the global minimum) when the first expression is smaller, giving the exact phase boundary
\begin{align}
\frac{I_2}{I_1} = \frac{p_1\!\left(1-\tfrac{3p_1}{4}\right)}{p_2\!\left(1-\tfrac{3p_2}{4}\right)}\,.
\end{align}

%or in terms of $\bar{I} = (I_1+I_2)/2$ and  $\delta{I} = (I_1-I_2)/2$,
%\begin{align}
%\frac{\partial E[\mathcal{L}]}{\partial\theta}\Bigg|_{b^*} &= 
%\frac{s_{2\theta}}{6}\Big[\bar{I}(p_1-p_2) + \delta I(p_1+p_2)\Big] \nonumber\\
%&\quad - \frac{p_1 p_2}{8}\Big[c_{2\theta}\left(\bar{I} - \delta I\, c_{2\theta}\right) + 2\delta I\, s_{2\theta}^2\Big] \nonumber\\
%&\quad - \frac{s_\theta(s_\theta-c_\theta)}{8}\bar{I}\Big[p_1^2(2s_{2\theta}-c_{2\theta}) - p_2^2\tfrac{c_\theta}{s_\theta}(2s_{2\theta}+c_{2\theta})\Big] \nonumber\\
%&\quad - \frac{s_\theta(s_\theta-c_\theta)}{8}\delta I\Big[p_1^2(2s_{2\theta}-c_{2\theta}) + p_2^2\tfrac{c_\theta}{s_\theta}(2s_{2\theta}+c_{2\theta})\Big]\,.
%\end{align}

\section{Non-Linear model}
\begin{align}
h &= w x\,, \qquad x' = \mathrm{ReLU}[w^T h + b] = \mathrm{ReLU}[w^T w\, x + b]\,.
\end{align}
We again take $n=2$, $m=1$, with the same parameterization
\begin{align}
w &= (c_\theta,\, s_\theta)\,, \qquad
w^T w = \begin{pmatrix} c^2_\theta & \tfrac{1}{2}s_{2\theta} \\ \tfrac{1}{2}s_{2\theta} & s^2_\theta \end{pmatrix}.
\end{align}
The network is trained on $\{x_{Ji}\}$ with $J=1,\ldots,k$ and $i=1,2$. Each $x_{Ji}$ is i.i.d.: zero with probability $1-p_i$, $\mathrm{Uniform}[0,1]$ with probability $p_i$. The loss function is a weighted MSE,
\begin{align}
\mathcal{L} = \frac{1}{2k}\sum_J\sum_i I_i\,(x_{Ji}-x'_{Ji})^2\,,
\end{align}
with moments $E[x_{Ji}]=p_i/2$ and $E[x_{Ji}^2]=p_i/3$.

For $k\gg 1$, $\mathcal{L}$ concentrates around its mean. We decompose $E[\mathcal{L}_J]$ according to which features are active, absorbing the probability weights into each term:
\begin{align}
E[\mathcal{L}_J] &= A + B + C + D\,.
\end{align}

\subsection{Both features inactive} ($x_1=x_2=0$, probability $(1-p_1)(1-p_2)$). When both inputs vanish, $x'_i = \mathrm{ReLU}(b_i)$, so
\begin{align}
A &= (1-p_1)(1-p_2)\sum_i I_i\,\mathrm{ReLU}(b_i)^2\,.
\end{align}
This piece is leading in the sparse-features regime and is independent of $\theta$; if both features are zero, the best reconstruction would for $b_i\leq0$. Thus, this term in the loss function pushes $b_i$ the non-positive values.
\subsection{Only feature 1 active} In this case, $x_1\sim\mathrm{Unif}[0,1]\,, x_2=0$, with probability $(1-p_2)p_1$. The reconstructions reduce to $x'_1=\mathrm{ReLU}\left(c_\theta^2 x_1+b_1\right)$ and $x'_2=\mathrm{ReLU}(\frac{s_{2\theta}}{2} x_1+b_2)$
We introduce the single-feature auxiliary integrals
\begin{align}
F_n(\alpha,\beta) &\equiv \int_0^1 \mathrm{ReLU}(\alpha x+\beta)^n\,\mathrm{d}x
= \frac{\mathrm{ReLU}(\alpha+\beta)^{n+1}-\mathrm{ReLU}(\beta)^{n+1}}{\alpha(n+1)}\,, \\[6pt]
G(\alpha,\beta) &\equiv \int_0^1 x\,\mathrm{ReLU}(\alpha x+\beta)\,\mathrm{d}x
= \frac{1}{\alpha}\!\left[\frac{\mathrm{ReLU}(\alpha+\beta)^3-\mathrm{ReLU}(\beta)^3}{3\alpha} - \beta\,F_1(\alpha,\beta)\right],
\end{align}
and find
\begin{align}
B &= (1-p_2)p_1\bigg[I_1\!\left(\tfrac{1}{3}-2G(c^2_\theta,b_1)+F_2(c^2_\theta,b_1)\right)
+ I_2\,F_2\!\left(\tfrac{s_{2\theta}}{2},b_2\right)\bigg]\,.
\end{align}
%We consider a few interesting cases. If $\theta=0$, this term becomes
%\begin{align}
%B(\theta=0) &= (1-p_2)p_1 \left[ I_2(\text{ReLU}\,b_2)^2+I_1\varepsilon (b_1) \right]\,,
%\\
%\varepsilon (x)  &\equiv \begin{cases}
%x^2 \;\;\;\;\;\;\;\;\; &0 \geq x
%\\
%x^2(1+x)-x^3/3  & -1\geq x<0
%\\
%1/3 & x<-1
%\end{cases}\,.
%\end{align}
%$\varepsilon$ is minimized at $x=0$, where $\varepsilon(0)=0$.
%On the other hand, if $x=\pi/2$,
%\begin{align}
%B(\theta=\pi/2) &= (1-p_2)p_1 \left[ I_2(\text{ReLU}\,b_2)^2+I_1\psi(b_1) \right]\,,
%\\
%\psi(x) &\equiv \begin{cases}
%1/3-x+x^2 \;\;\;\;\;\;\;\;\; &0 \geq x
%\\
%1/3 & x<0
%\end{cases}\,.
%\end{align}
%$\psi$ is minimized at $x=1/2$, where $\psi(1/2) = 1/12$.
\subsection{Only feature 2 active} 
By symmetry ($1\leftrightarrow 2$, $c_\theta\leftrightarrow s_\theta$):
\begin{align}
C &= (1-p_1)p_2\bigg[I_1\,F_2\!\left(\tfrac{s_{2\theta}}{2},b_1\right)
+ I_2\!\left(\tfrac{1}{3} - 2G(s^2_\theta,b_2) + F_2(s^2_\theta,b_2)\right)\bigg]\,.
\end{align}
\subsection{No superposition}
Consider the $\theta=0$ phase. The loss function is minimized with the choice $b_1=0$, end one ends up with
\begin{align}
(A+B+C)\bigg |_{\theta=0,b_1=0} = I_2 \left[ \text{ReLU}^2(b_2)+p_2\left( \frac13 -\text{ReLU}(b_2)\right)\right] +\mathcal{O}(p^2)
\end{align}
The loss is then minimzed for $b_2=p_2/2$, and we find
\begin{align}
(A+B+C)\bigg |_{\theta=0,b_1=0,b_2=p_2/2} = \frac{I_2 p_2}{3}\,. 
\end{align}
By symmetry
\begin{align}
(A+B+C)\bigg |_{\theta=\pi/2,b_2=0,b_1=p_1/2} = \frac{I_1 p_1}{3}\,. 
\end{align}
\subsection{Superposition}
Let us consider $\theta=3\pi/4$. In this case, the calculation is more involved since different values of $b_1,b_2$ lead to different ReLU's being activated. Nevertheless, after a straight forward calculation, one finds that the loss function is minimized for $b_1,b_2 \to 0^+$, and 
\begin{align}
(A+B+C)\bigg |_{\theta=3\pi/4,b_1=b_2=0} = \frac{I_1 p_1+I_2 p_2}{12}-\mathcal{O}(p^2)\,.
\end{align}
which allows us to determine the phase structure 
\begin{align}
\theta = \begin{cases}
0 & I_2 p_2<(I_1 p_1)/3
\\
3\pi/4\;\;\;\;\; &(I_1 p_1)/3< I_2 p_2 < 3 I_1 p_1
\\
\pi/2 & 3 I_1 p_1<I_2 p_2
\end{cases}\,.
\end{align}

%Also from symmetry,
%\begin{align}
%C(\theta=0) &= B(\theta=\pi/2,1 \leftrightarrow2) =  (1-p_1)p_2 \left[ I_1(\text{ReLU}\,b_1)^2+I_2\psi(b_2) \right]\,,
%\\
%C(\theta=\pi/2) &= B(\theta=0,1 \leftrightarrow2) =  (1-p_1)p_2 \left[ I_1(\text{ReLU}\,b_1)^2+I_2\varepsilon (b_2) \right]\,.
%\end{align}
%Putting it all together, and neglecting $p_1 p_2$,
%\begin{align}
%(B+C)\bigg |_{\theta=0} = p_1 \left[ I_2(\text{ReLU}\,b_2)^2+I_1\varepsilon (b_1) \right]+p_2 \left[ I_1(\text{ReLU}\,b_1)^2+I_2\psi(b_2) \right]
%\end{align}
%$b_1=0$ minimizes the contribution from $b_1$; leaving us with
%\begin{align}
%(A+B+C)\bigg |_{\theta=0,b_1=0} &= (1-p_1-p_2)I_2\text{ReLU}(b_2)^2  + I_2 p_2  \left[ (p_1/p_2)(\text{ReLU}\,b_2)^2 + \psi(b_2) \right]
%\\
%&\approx I_2 (1-p_2)  \left[ \text{ReLU}(b_2)^2  +   p_2 \psi(b_2)  \right]
%\end{align}
%The minimum is found where $b_2 = p_2/2$,
%\begin{align}
%(A+B+C)\bigg |_{\theta=0,b_1=0,b_2=p_2/2} = \frac{I_2 p_2}{3}
%\end{align}
%And from symmetry,
%\begin{align}
%(A+B+C)\bigg |_{\theta=\pi/2,b_2=0,b_1=p_1/2} = \frac{I_1 p_1}{3}\,.
%\end{align}
%Now let's check what happens at $\pi/4$. In this case,


\newpage
\textbf{Both features active} ($x_1,x_2\sim\mathrm{Unif}[0,1]$, probability $p_1 p_2$). Now both inputs enter the reconstruction, coupling $x_1$ and $x_2$. Integrating over both requires the double auxiliary integrals
\begin{align}
\mathcal{F}_n(\alpha,\gamma,\beta) &\equiv \int_0^1 F_n(\alpha,\gamma y+\beta)\,\mathrm{d}y \nonumber\\
&= \frac{\mathrm{ReLU}(\alpha+\gamma+\beta)^{n+2}-\mathrm{ReLU}(\gamma+\beta)^{n+2}-\mathrm{ReLU}(\alpha+\beta)^{n+2}+\mathrm{ReLU}(\beta)^{n+2}}{\alpha\gamma(n+1)(n+2)}\,, \\[8pt]
\mathcal{H}(\alpha,\gamma,\beta) &\equiv \int_0^1 y\,F_1(\alpha,\gamma y+\beta)\,\mathrm{d}y \nonumber\\
&= \frac{\mathrm{ReLU}(\alpha+\gamma+\beta)^3-\mathrm{ReLU}(\gamma+\beta)^3}{6\alpha\gamma}
- \frac{\mathrm{ReLU}(\alpha+\gamma+\beta)^4-\mathrm{ReLU}(\alpha+\beta)^4-\mathrm{ReLU}(\gamma+\beta)^4+\mathrm{ReLU}(\beta)^4}{24\alpha\gamma^2}\,, \\[8pt]
\mathcal{G}(\alpha,\gamma,\beta) &\equiv \int_0^1 G(\alpha,\gamma y+\beta)\,\mathrm{d}y
= \frac{1}{\alpha}\Big[\mathcal{F}_2(\alpha,\gamma,\beta) - \beta\,\mathcal{F}_1(\alpha,\gamma,\beta) - \gamma\,\mathcal{H}(\alpha,\gamma,\beta)\Big]\,,
\end{align}
giving
\begin{align}
D &= p_1 p_2\bigg[I_1\!\left(\tfrac{1}{3}-2\mathcal{G}\!\left(c^2_\theta,\tfrac{s_{2\theta}}{2},b_1\right)+\mathcal{F}_2\!\left(c^2_\theta,\tfrac{s_{2\theta}}{2},b_1\right)\right) \nonumber\\
&\qquad\quad + I_2\!\left(\tfrac{1}{3}-2\mathcal{G}\!\left(s^2_\theta,\tfrac{s_{2\theta}}{2},b_2\right)+\mathcal{F}_2\!\left(s^2_\theta,\tfrac{s_{2\theta}}{2},b_2\right)\right)\bigg]\,.
\end{align}

Combining all four contributions, the expected loss is
\begin{align}
\boxed{E[\mathcal{L}] = \frac{1}{2}(A+B+C+D)\,.}
\end{align}

\begin{thebibliography}{9}
\end{thebibliography}
\end{document}